% Related Work
\section{Related Work}

\subsection{Advantage-Weighted Regression (AWR) and Related Methods}

Advantage-Weighted Regression (AWR) \citep{peng2019advantage} introduced exponentially-weighted supervised learning for offline reinforcement learning in robotics:
\begin{equation}
    \mathcal{L}_{\text{AWR}} = \mathbb{E}\left[ \exp\left(\frac{A(s,a)}{\beta}\right) \cdot \log \pi_\theta(a|s) \right]
\end{equation}

AWR computes advantages using TD($\lambda$) returns from a separately trained value function. Advantage-weighted Policy Actor (APA) extends this with squared error stabilization for online learning.

\textbf{Critical distinction}: WMTP does \textit{not} reproduce AWR in its entirety. We adopt only the \textbf{$\exp(A/\beta)$ weighting form}, while making significant departures in advantage computation (GAE instead of TD($\lambda$) returns) and value learning (pairwise ranking instead of MSE regression). Table~\ref{tab:awr_comparison} summarizes these differences.

\begin{table}[h]
\centering
\caption{Comparison of AWR and WMTP components. WMTP adopts only the exponential weighting form from AWR while using different methods for advantage computation and value learning.}
\label{tab:awr_comparison}
\begin{tabular}{lcc}
\toprule
\textbf{Component} & \textbf{AWR} \citep{peng2019advantage} & \textbf{WMTP (Ours)} \\
\midrule
Domain & Robotics control & LLM code generation \\
Advantage computation & TD($\lambda$) return-based & \textbf{GAE (TD-error accumulation)} \\
Value learning & Regression (MSE) & \textbf{Pairwise Ranking} \\
Weight form & $\exp(A/\beta)$ & $\exp(A/\beta)$ (same form only) \\
\bottomrule
\end{tabular}
\end{table}

\subsection{Implicit Q-Learning and Offline RL for LLMs}

Several works apply offline RL principles to LLM fine-tuning:

\begin{description}[style=unboxed,leftmargin=0cm]
    \item[Q-SFT:] Interprets token probabilities as Q-values and uses Bellman targets for weighting. Requires careful tuning to avoid distributional shift.

    \item[IQL/ILQL \citep{snell2022offline}:] Implicit Language Q-Learning uses expectile regression to learn conservative value estimates. However, it requires \textbf{inference-time logit perturbation}, adding computational overhead during generation.

    \item[OREO:] Uses policy gradient and odds ratio formulations, which can be unstable due to importance sampling variance.
\end{description}

WMTP occupies a distinct position in this landscape:
\begin{itemize}[nosep]
    \item \textbf{Weighted SFT}: Stable supervised learning (no policy gradient variance)
    \item \textbf{Training-time weighting}: Zero inference overhead (critic not used at test time)
    \item \textbf{Explicit value}: Interpretable token-level scores for debugging and analysis
\end{itemize}

\subsection{Multi-Token Prediction}

Multi-Token Prediction (MTP) \citep{gloeckle2024better} extends standard next-token prediction by using multiple independent heads to predict $K$ future tokens simultaneously:
\begin{equation}
    \mathcal{L}_{\text{MTP}} = \frac{1}{K} \sum_{k=1}^{K} \sum_t \text{CE}(z_{t,k}, y_{t+k})
\end{equation}

MTP improves sample efficiency and representation quality through multi-step prediction. However, the loss averages uniformly across all tokens and heads, without distinguishing token importance.

\textbf{WMTP contribution}: We maintain the MTP multi-head architecture while \textbf{redistributing learning signal based on token-level value estimates}. This allows the model to focus on tokens that matter most for solution correctness.

\subsection{Summary: Method Positioning}

Table~\ref{tab:method_comparison} positions WMTP relative to existing approaches:

\begin{table}[h]
\centering
\caption{Comparison of related methods. WMTP combines weighted SFT stability with training-time-only critic usage and MTP support.}
\label{tab:method_comparison}
\resizebox{\textwidth}{!}{%
\begin{tabular}{lcccccc}
\toprule
\textbf{Method} & \textbf{Loss} & \textbf{Weighting} & \textbf{Value Learning} & \textbf{Extra Model} & \textbf{MTP Support} & \textbf{Inference Overhead} \\
\midrule
Standard MTP & CE & Uniform & -- & -- & Yes & None \\
AWR & Weighted CE & $\exp(A/\beta)$ & Regression & Critic & No & None \\
ILQL & Modified CE & Q-based & Implicit & -- & No & \textbf{Yes} \\
\textbf{WMTP (Ours)} & \textbf{Weighted CE} & $\exp(A/\beta)$ & \textbf{Pairwise} & \textbf{Critic} & \textbf{Yes} & \textbf{None} \\
\bottomrule
\end{tabular}%
}
\end{table}
